% Computer Science A --- Spring Boot (\texttt{application.*}, auto-config, actuator, packaging, Boot testing). See \texttt{spring.tex} for the core framework.

\runningheader{Computer Science A}{Spring Boot}{Page \thepage\ of \numpages}

\begingradingrange{csaspringboot}

\question[4]
Overall project configuration can be provided to a Spring application using which files?
\begin{choices}
  \choice \texttt{application.xml} and \texttt{application.yml}
  \choice \texttt{application.properties} and \texttt{application.xml}
  \choice \texttt{application.xml} and \texttt{bootstrap.xml}
  \CorrectChoice \texttt{application.yml} and \texttt{application.properties}
\end{choices}
\begin{solution}
  Spring Boot convention loads \textbf{\texttt{application.properties}} or \textbf{\texttt{application.yml}} (and profile-specific variants) from the classpath. \texttt{bootstrap} files belong to Spring Cloud Context, not the baseline pair in the stem.
\end{solution}

\question[4]
Using the default configuration, Spring Boot includes an embedded
\begin{choices}
  \CorrectChoice Tomcat server
  \choice Postgres server
  \choice in-memory database
  \choice user
\end{choices}
\begin{solution}
  The default embedded servlet container is \textbf{Apache Tomcat}. You can switch to Jetty or Undertow via starters. An in-memory database is optional via dependencies (e.g.\ H2), not the default embedded HTTP stack.
\end{solution}

\question[4]
Spring Boot features an ``opinionated configuration'' which means that:
\begin{choices}
  \choice It requires you to tell it how to configure all functionality.
  \CorrectChoice When not provided with configuration, it uses reasonable default configurations.
  \choice It cannot be configured.
  \choice You will never need to configure it.
\end{choices}
\begin{solution}
  \textbf{Opinionated defaults} (starters, auto-configuration) get a project running quickly; you override only what you need. You can still configure extensively---it is not ``unconfigurable.''
\end{solution}

\question[4]
Spring Boot actuators can be used to
\begin{choices}
  \choice interact with your server using a convenient GUI.
  \choice open a running application on the command line.
  \CorrectChoice get runtime information about the application via HTTP.
  \choice start a server when it is not running.
\end{choices}
\begin{solution}
  \textbf{Actuators} expose operational endpoints (health, metrics, env, \ldots), often over HTTP (and JMX), for monitoring and management. They are not a general GUI, a shell, or a server launcher.
\end{solution}

\question[4]
The Spring Boot actuator automatically restarts your Spring application whenever files on the classpath change.
\begin{choices}
  \choice TRUE
  \CorrectChoice FALSE
\end{choices}
\begin{solution}
  Automatic restart on classpath changes is provided by \textbf{Spring Boot DevTools}, not by Actuator alone. Actuator focuses on observability and management endpoints.
\end{solution}

\question[4]
By default all the actuator endpoints are enabled except for \texttt{shutdown}.
\begin{choices}
  \CorrectChoice TRUE
  \choice FALSE
\end{choices}
\begin{solution}
  This follows the classic teaching summary: many endpoints are exposed by default in older defaults; \texttt{shutdown} is sensitive and typically off unless explicitly enabled. Exact exposure changed across Boot 2.x/3.x (e.g.\ only \texttt{health} over HTTP unless configured); treat this as the keyed conceptual answer and align with your Boot version in lecture.
\end{solution}

\question[4]
What kind of file will a Spring Boot project typically be packaged as?
\begin{choices}
  \choice EXE file
  \choice BAT file
  \choice CMD file
  \CorrectChoice JAR file
\end{choices}
\begin{solution}
  Executable \textbf{fat JAR}s (\texttt{spring-boot-maven-plugin} / Gradle plugin) bundle dependencies and \texttt{Main-Class} for \texttt{java -jar}. WAR deployment to an external servlet container is also common but the stem targets the typical standalone artifact.
\end{solution}

\question[4]
What is Spring Boot autoconfiguration?
\begin{choices}
  \choice Autoconfiguration automatically wires up your bean dependencies and adds them to the IOC container
  \CorrectChoice Autoconfiguration attempts to automatically configure your Spring application based on the jar dependencies that you have added
  \choice Autoconfiguration sets up REST API endpoints based on the classes configured in your application
  \choice Autoconfiguration sets the Java version that your Spring project will use based on the features that your code uses
\end{choices}
\begin{solution}
  \textbf{@EnableAutoConfiguration} (via \texttt{@SpringBootApplication}) pulls in conditional configuration classes on the classpath (e.g.\ DataSource when JDBC starter + properties present). DI wiring is broader than auto-config alone; REST endpoints come from your controllers, not auto-config wholesale.
\end{solution}

\question[4]
Spring Boot Actuator custom endpoints can be created by\ldots
\begin{choices}
  \choice Writing MockMVC integration tests with custom data
  \choice Writing a controller that uses the \texttt{@RequestMapping} annotation
  \choice Putting the custom endpoint path in the \texttt{application.yml} file and linking to the implementing class
  \CorrectChoice Annotating a bean with the \texttt{@Endpoint} annotation and give an \texttt{id} property to specify the endpoint path
\end{choices}
\begin{solution}
  Custom \textbf{Actuator} endpoints (JMX or HTTP, depending on exposure settings) use the actuator programming model (\texttt{@Endpoint}, \texttt{@ReadOperation}, etc.). A normal \texttt{@RestController} serves application APIs, not actuator infrastructure.
\end{solution}

\question[4]
How can you test asynchronous code in Spring Boot?
\begin{choices}
  \choice Use the \texttt{@Async} annotation
  \choice Use the \texttt{@TestAsync} annotation
  \CorrectChoice Use the \texttt{CompletableFuture} class
  \choice Asynchronous testing is not supported in Spring Boot
\end{choices}
\begin{solution}
  The keyed answer follows the source. In practice, \texttt{@Async} enables async execution; testing often uses \texttt{CompletableFuture.get()}, \texttt{@Test} timeouts, or Spring's async test utilities (\texttt{WebTestClient}, Awaitility, etc.). There is no standard \texttt{@TestAsync} in core Spring Test.
\end{solution}

\question[4]
How can you customize the test environment using profiles in Spring Boot?
\begin{choices}
  \CorrectChoice Use the \texttt{@ActiveProfiles} annotation
  \choice Use the \texttt{@TestConfiguration} annotation
  \choice Use the \texttt{application-test.properties} file
  \choice Profiles cannot be used to customize the test environment
\end{choices}
\begin{solution}
  \texttt{@ActiveProfiles("test")} on a test class activates profile-specific beans and property files (e.g.\ \texttt{application-test.yml}). \texttt{@TestConfiguration} supplies extra beans; \texttt{application-test.properties} is common but the annotation directly drives which profiles load in tests.
\end{solution}

\question[4]
What is the primary purpose of Spring Boot Testing?
\begin{choices}
  \choice To write production code
  \CorrectChoice To test individual units of code
  \choice To design the application architecture
  \choice To deploy the application in production
\end{choices}
\begin{solution}
  The keyed answer follows the source (unit-focused wording). Spring Boot's test stack also supports \textbf{slice} tests (\texttt{@WebMvcTest}, \texttt{@DataJpaTest}) and full \textbf{integration} tests (\texttt{@SpringBootTest}); clarify in class that ``testing'' spans multiple levels.
\end{solution}

\question[4]
Which testing framework is integrated by default in Spring Boot for writing tests?
\begin{choices}
  \CorrectChoice JUnit
  \choice TestNG
  \choice Mockito
  \choice Cucumber
\end{choices}
\begin{solution}
  \texttt{spring-boot-starter-test} brings in \textbf{JUnit Jupiter} (JUnit 5), AssertJ, Mockito, Hamcrest, \texttt{spring-test}, etc. TestNG and Cucumber are optional additions.
\end{solution}

\question[4]
Explain the role of the \texttt{@WebMvcTest} annotation in Spring Boot Testing.
\begin{choices}
  \CorrectChoice It is used for testing web layers only
  \choice It is used for testing database interactions
  \choice It is used for testing service layers only
  \choice It is used for testing repositories
\end{choices}
\begin{solution}
  \texttt{@WebMvcTest} slices the context to \textbf{web layer} components (controllers, JSON converters, \texttt{MockMvc}) and excludes full data layer auto-configuration unless you import extras. Full-stack tests use \texttt{@SpringBootTest}.
\end{solution}

\question[4]
Explain the purpose of the \texttt{@AutoConfigureMockMvc} annotation in Spring Boot Integration Testing.
\begin{choices}
  \choice It configures the application to use a mock database
  \CorrectChoice It auto-configures the \texttt{MockMvc} instance for testing
  \choice It configures the application to use a real database
  \choice It auto-configures the application context for unit testing
\end{choices}
\begin{solution}
  \texttt{@AutoConfigureMockMvc} registers \texttt{MockMvc} for sending fake HTTP requests in tests (often with \texttt{@SpringBootTest}). It does not imply mock vs.\ real DB; that is controlled by test slices, \texttt{@MockBean}, \texttt{@DataJpaTest}, or Testcontainers.
\end{solution}

\uplevel{\par\textbf{Spring Boot practice bank}\par}

\question[4]
Which annotation usually serves as the primary entry-point annotation for a Spring Boot app?
\begin{choices}
  \CorrectChoice \texttt{@SpringBootApplication}
  \choice \texttt{@EnableJpaRepositories}
  \choice \texttt{@EnableScheduling}
  \choice \texttt{@ConfigurationPropertiesScan} only
\end{choices}
\begin{solution}
  \texttt{@SpringBootApplication} combines \texttt{@Configuration}, \texttt{@EnableAutoConfiguration}, and component scanning, making it the common application bootstrap annotation.
\end{solution}

\question[4]
How does Spring Boot usually decide whether to auto-configure a bean?
\begin{choices}
  \choice It always creates every possible bean from \texttt{spring-boot-autoconfigure}
  \CorrectChoice It uses conditions (classpath, existing beans, properties, environment) to activate matching auto-configuration classes
  \choice It only reads XML files from \texttt{/config}
  \choice It requires manual \texttt{new} calls in the main method
\end{choices}
\begin{solution}
  Auto-configuration is conditional, commonly via \texttt{@ConditionalOnClass}, \texttt{@ConditionalOnMissingBean}, and property-based conditions.
\end{solution}

\question[4]
Which property best sets the HTTP server port to 9090 in Spring Boot?
\begin{choices}
  \CorrectChoice \texttt{server.port=9090}
  \choice \texttt{spring.http.port=9090}
  \choice \texttt{tomcat.server.port=9090}
  \choice \texttt{boot.port=9090}
\end{choices}
\begin{solution}
  \texttt{server.port} is the standard Boot server property in \texttt{application.properties} (or YAML equivalent).
\end{solution}

\question[4]
In current Boot defaults, which Actuator endpoint is typically exposed over HTTP out of the box?
\begin{choices}
  \CorrectChoice \texttt{/actuator/health}
  \choice \texttt{/actuator/shutdown}
  \choice \texttt{/actuator/env}
  \choice All endpoints are exposed by default
\end{choices}
\begin{solution}
  Modern Boot versions generally expose only a minimal set over HTTP by default (commonly health). Additional endpoints require explicit exposure configuration.
\end{solution}

\question[4]
When should you prefer \texttt{@SpringBootTest} over \texttt{@WebMvcTest}?
\begin{choices}
  \choice When you only want to test controller request mapping and JSON serialization in isolation
  \CorrectChoice When you need broader integration coverage across the full application context
  \choice When you want to exclude service and repository beans
  \choice Never; \texttt{@WebMvcTest} always fully replaces it
\end{choices}
\begin{solution}
  \texttt{@WebMvcTest} is a web slice; \texttt{@SpringBootTest} loads a wider context and is used for integration-style scenarios spanning multiple layers.
\end{solution}

\question[4]
Which statement about externalized configuration precedence is generally correct in Spring Boot?
\begin{choices}
  \choice Values in packaged \texttt{application.properties} always override command-line arguments
  \CorrectChoice Command-line arguments typically have high precedence and can override packaged config values
  \choice Environment variables are ignored if a properties file exists
  \choice Profile-specific files are never loaded
\end{choices}
\begin{solution}
  Boot combines multiple property sources with precedence rules; command-line options are high-priority and commonly override defaults from files.
\end{solution}

\endgradingrange{csaspringboot}
